\begin{frame}{Source identities must be recorded before counting separate evidence.\ev{\tagO}}
\begin{tabularx}{\linewidth}{L{3.4cm}L{3.4cm}Y}\toprule
Visible count & Relationship to retain & Claim boundary \\\midrule
Cells & Patient identity & Cells may share one source \\
Mining addresses & Linked agents & Addresses can belong to one actor \\
Executions & Observation and case IDs & New execution can reuse old evidence \\\bottomrule\end{tabularx}
\claim{Record the unit of analysis and its relationships with each result.}
\sources{Blackburn et al., \href{https://arxiv.org/abs/2206.02871}{arXiv:2206.02871}; Saurty-Seerunghen et al., iScience 2026; provenance semirings, PODS 2007.}
\qadetail{The coauthored Bitcoin study linked most early mining to sixty-four agents; malignant-cell clustering depended on patient identity. These required domain-specific inference, not a generic lineage algorithm. They do not measure fork outcomes. Proposed runtime manifest: candidate digest, parent snapshot, input tree, execution ID, observation IDs, prompt/model version, seed policy, case/fixture IDs and feedback IDs. Built lineage transport does not estimate effects or hidden shared blind spots. Additional perception/physical-network and ENCODE examples remain methodological context in REFERENCES.md.}
\note{[0:50] Across our coauthored studies, the visible identifier was sometimes finer than the source of the observation. Tumor cells shared patients. Mining addresses could be linked to the same agent. In a runtime, executions can share a test case or return the same observation. The recurring step is to identify the unit and record its relationships before counting it. These domain examples do not estimate fork correlation. They motivate the manifest, while the next question concerns the population from which a checkpoint observation is drawn.}
\end{frame}
